% ===================================================================
%  TAULA N° 1 — L'Escòla. — L'Institut. — La Classa.
%
%  Ceci n'est PAS une transcription : c'est une traduction, faite en
%  2026, en occitan languedocien d'aujourd'hui, graphie classique.
%  D'ou trois differences avec texto/io et texto/fr, et trois
%  seulement :
%
%   * pas de \nl ni de \cc — la traduction ne suit aucune ligne
%     imprimee, et les lignes de ce fichier ne sont que du confort ;
%   * pas de \begin{VUpage} — elle n'a ni feuillet ni folio, et la
%     page de lecture ne lui en affiche pas ;
%   * le gras se pose sur le TERME seul, ponctuation dehors.
%
%  Tout le reste est identique, et doit l'etre : les cles « %%K » sont
%  celles de texto/io, macro pour macro pour l'apparat, et les renvois
%  \textsuperscript{(n)} visent les MEMES objets de la planche — ce
%  sont eux qui ouvrent les gros plans.
%
%  LA SEPTIEME DES DIX-SEPT LANGUES IDISTES. L'ido a eu ses groupes
%  dans le Midi des les annees 1910 ; et l'occitan a de plus, dans
%  cette histoire, un titre particulier : c'est de sa parente que le
%  catalan tient la sienne, et les deux ensemble donnent a la serie
%  deux temoins du meme rameau.
%
%  LA COLONNE SE TRADUIT SUR L'IDO ET SE CONTROLE SUR texto/ca, qui
%  est la plus proche de toutes. Mais elle ne s'y recopie pas : les
%  deux langues ont diverge depuis huit siecles, et le vocabulaire
%  scolaire n'est pas le meme. « aula » se dit « classa »,
%  « pissarra » se dit « tablèu negre », « faristol » se dit
%  « pupitre », « esbarjo » se dit « recreacion ».
%
%  MEME ORDRE DES MOTS QU'EN CATALAN, donc meme ordre des renvois :
%  le complement du nom se postpose, et l'ido s'y ecrit sans qu'un
%  seul renvoi ait a bouger.
%
%  LES NOMS DE PERSONNE SONT OCCITANS : Ioannes est Joan, Iakobus est
%  Jaume, Karolus est Carles. Seuls le chien Medor et l'ane Aliboron
%  gardent leur nom. Les noms latins de la note (*) ne bougent pas :
%  c'est d'eux que la note parle.
% ===================================================================

%%K t01-apar-0 apar
\VUsaut{2.0mm}
\VUcentre{12.6pt}{120}{FULHET EXPLICATIU}
\VUsaut{3.2mm}
\VUcentre{8.4pt}{160}{DE}
\VUsaut{2.6mm}
\VUtitre{15.6pt}{0}{las Taulas Auxiliàrias Delmas}
\VUsaut{3.4mm}
\VUornamento{9pt}
\VUsaut{5.0mm}
\VUcentre{13.2pt}{90}{PRIMIÈRA SÈRIA}
\VUsaut{3.0mm}
\VUfilet{16mm}
\VUsaut{5.6mm}

%%K t01-apar-1 apar
\VUtitre{15.0pt}{60}{TAULA N° 1}
\VUsaut{1.4mm}
\VUtitre{11.4pt}{0}{L'Escòla. --- L'Institut. --- La Classa.}
\VUsaut{2.2mm}
\VUfilet{20mm}
\VUsaut{6.4mm}

%%K t01-tit-1 sub
\VUpk{10.2pt}{40}{Descripcion generala.}
\VUsaut{3.0mm}

%%K t01-01-1 p
1. --- Aquesta taula representa una \VUgras{classa} d'\VUgras{institut}. Dins
aquesta classa i a uèch \VUgras{taulas} \textsuperscript{(18)} e uèch
\VUgras{bancs} \textsuperscript{(32)}. Sus cada banc s'assètan tres escolans. Lo
\VUgras{professor} \textsuperscript{(7)} es assetat sus una \VUgras{cadièira}
\textsuperscript{(8)} a sa \VUgras{catedra} \textsuperscript{(6)}.

%%K t01-02-1 p
2. --- A esquèrra de la catedra i a un \VUgras{armari} amb pòrtas de veire
\textsuperscript{(15)}. A drecha i a lo \VUgras{tablèu negre}
\textsuperscript{(1)} amb lo \VUgras{netejador} \textsuperscript{(2)} e lo
\VUgras{gis} \textsuperscript{(3)}.

%%K t01-02-2 p
Dessús la catedra pendolan de \VUgras{laminas muralas}
\textsuperscript{(9, 11, 12)} e una \VUgras{mapa} \textsuperscript{(10)}.

%%K t01-03-1 p
3. --- Totes los escolans son pas a sa \VUgras{plaça} \textsuperscript{(17)}. Joan
\textsuperscript{(4)} es drech contra lo tablèu; ten lo netejador e escafa las
\VUgras{figuras de geometria}.

%%K t01-03-2 p
Pèire es a prèp de la catedra; recampa los \VUgras{exercicis escriches}
\textsuperscript{(5)} e los pòrta al professor.

%%K t01-04-1 p
4. --- \VUgras{Aqueste} professor es lo professor de \VUgras{lengas modèrnas}.
Ensenha als escolans de parlar, legir e escriure de lengas estrangièras
\VUgras{(alemand, anglés, espanhòl, italian, rus} o \VUgras{francés)}.

%%K t01-05-1 p
5. --- La classa es fòrça vasta e fòrça clara. Al fons i a una \VUgras{fenèstra}
larga amb dos \VUgras{montants} \textsuperscript{(78)} e una \VUgras{pòrta
veirada} per l'aerejar e l'esclairar.

%%K t01-05-2 p
Aquesta classa es pas freja. Al mitan, per l'escalfar, i a una \VUgras{estufa}
\textsuperscript{(46)} fòrça bona amb son \VUgras{tuble} \textsuperscript{(45)}.

%%K t01-05-3 p
Sus la paret i a de \VUgras{penjadors} \textsuperscript{(58)} fixats; ne pendolan
las casquetas e los mantèls dels escolans.

%%K t01-tit-2 sub
\VUsaut{5.2mm}
\VUpk{10.2pt}{40}{La cort de recreacion.}
\VUsaut{3.6mm}

%%K t01-06-1 p
6. --- Lo \VUgras{relòtge} \textsuperscript{(63)} marca nòu oras e cinc.

%%K t01-06-2 p
La \VUgras{recreacion} \textsuperscript{(76)} ven d'acabar e la classa tòrna
començar. La recreacion se fa dins la \VUgras{cort} \textsuperscript{(75)}. Vesèm
qualques escolans que jògan. Un \VUgras{mèstre d'estudi} \textsuperscript{(74)}
los susvelha.

%%K t01-07-1 p
7. --- Dins la cort vesèm en mai d'autras personas: l'\VUgras{inspector}
\textsuperscript{(73)}, lo \VUgras{director} \textsuperscript{(71)} e lo
\VUgras{cap d'estudis} \textsuperscript{(72)}.

%%K t01-07-2 p
L'inspector inspècta las escòlas, lo director dirigís l'institut, lo cap
d'estudis susvelha los estudis.

%%K t01-07-3 p
Lo quatren personatge es lo \VUgras{concierge} \textsuperscript{(77)}. Pòrta jos
lo braç un registre per i marcar los absents.

%%K t01-08-1 p
8. --- Al fons de la cort i a los \VUgras{aparelhs de gimnastica}
\textsuperscript{(70)} e l'obrador de \VUgras{trabalhs manuals}
\textsuperscript{(69)}.

%%K t01-08-2 p
A drecha de la taula començam de veire las \VUgras{salas} de \VUgras{musica} e de
\VUgras{dessenh}.

%%K t01-noto-1 noto
\VUnotes{143.2mm}{%
(*) Pels \textit{noms de batisme} se conselha de los transcriure tanben dins sa
forma latina. Es l'unic mejan d'evitar de variantas sens nombre. D'alhors, cossí
causir entre \textit{Giovanni, Jean,} \textit{Johann, Jan, Hans, John, Ivan,}
etc.? Caldrà crear un diccionari especial dels noms de batisme? Prenguem del
latin son nominatiu e disèm: \VUgras{Ioannes, Ioanna; Iulius, Iulia; Iakobus;
Andreas; Lukas.} (Gramatica completa).%
}

%%K t01-tit-3 sub
\VUpk{10.2pt}{40}{Descripcion detalhada.}
\VUsaut{2.4mm}

%%K t01-tit-4 sub
\VUpk{10.2pt}{40}{Los bons escolans.}
\VUsaut{3.0mm}

%%K t01-09-1 p
9. --- Los escolans del primièr banc a drecha de la catedra son \VUgras{Enric}
(*), \VUgras{Estève} e \VUgras{Carles}.

%%K t01-09-2 p
Enric es fòrça atentiu e trabalhador; escota lo professor. Sus sa taula a un
\VUgras{quasèrn} \textsuperscript{(28)}, un \VUgras{libre} \textsuperscript{(29)},
un \VUgras{diccionari} \textsuperscript{(30)} e un \VUgras{plumièr}
\textsuperscript{(31)}. Son libre a fòrça \VUgras{paginas}; cada capítol conten
mantun \VUgras{paragraf} e cada \VUgras{linha} fòrça \VUgras{mots}.

%%K t01-09-4 p
Son vesin Estève \textsuperscript{(24)} es aplicat. Marca sus son quasèrn la
leiçon per lo còp que ven. A una polida \VUgras{escritura}. Son libre es barrat.
Vesèm pas que la \VUgras{cobèrta} \textsuperscript{(27)}. A son costat i a son
\VUgras{compàs} \textsuperscript{(26)}.

%%K t01-09-5 p
Carles \textsuperscript{(23)} ten son libre a la man; lo dobrís per repassar la
leiçon. Coma totes los escolans, a davant el un \VUgras{tintièr}
\textsuperscript{(25)}. Es obedient.

%%K t01-10-1 p
10. --- A la taula del costat i a \VUgras{Loís, Antòni} e \VUgras{Pau}. Loís
recita sa \VUgras{leiçon} \textsuperscript{(22)} e respond a las
\VUgras{questions} del professor. Antòni \textsuperscript{(21)} es fòrça
destorbat; de còps lo professor lo castiga quitament. Pau
\textsuperscript{(20)} es bon companh, amable e servicial. Es fòrça atentiu.

%%K t01-tit-5 sub
\VUsaut{4.2mm}
\VUpk{10.2pt}{40}{Los marrits escolans.}
\VUsaut{3.2mm}

%%K t01-11-1 p
11. --- Darrièr Enric i a \VUgras{Jòrdi} \textsuperscript{(38)}. Es pas un marrit
dròlle, mas es \VUgras{charraire}, destorbat e bolegadís. Sa
\VUgras{leugieretat} e sa \VUgras{desobeïssença} causan de \VUgras{desòrdre}
pendent la classa, e sovent merita una severa \VUgras{avertida}. Son
\VUgras{quasèrn de borrolhon} \textsuperscript{(35)} es brut, l'emplena de
\VUgras{tacas de tinta} amb sa \VUgras{pluma} \textsuperscript{(34)}, e per aquò
se sèrv totjorn de sa \VUgras{goma d'escafar} \textsuperscript{(36)}; sa
\VUgras{cartabèla} \textsuperscript{(33)} es esquiçada, que es fòrça descurós.
Perqué pòrta son \VUgras{portalapis} \textsuperscript{(37)} a la classa de lengas
modèrnas?

%%K t01-11-2 p
Son vesin Edmond \textsuperscript{(39)} l'aima pas. Es vertat qu'es un pauc
pigre, mas es calat e se ten quiet. Charra pas; solament que, al luòc
d'escotar, aima mai legir los \VUgras{contes} de son \VUgras{libre de lectura}.

%%K t01-11-3 p
Lo tresen escolan d'aqueste banc es \VUgras{Loís} \textsuperscript{(40)}. Es
aplicat. Escriu sus un \VUgras{fuèlh de papièr} blanc. Davant el a daissat son
\VUgras{papièr secant} \textsuperscript{(41)}.

%%K t01-12-1 p
12. --- Los escolans del segond banc, a esquèrra del professor, son fòrça joves.
Lo primièr, lo pichon \VUgras{Emili}, sap pas encara escriure plan; pòrta sa
\VUgras{lausa} \textsuperscript{(42)}, son \VUgras{lapis de lausa} e sa
\VUgras{esponga} \textsuperscript{(43)}.

%%K t01-12-2 p
A son costat i a \VUgras{Auguste} e \VUgras{Bernat} \textsuperscript{(44)}.
Aqueste darrièr trabalha plan, mas son \VUgras{aspècte} es fòrça descurós.

%%K t01-13-1 p
13. --- Darrièr eles i a tres \VUgras{pigres}. \VUgras{Albèrt} apren pas sas
leiçons. \VUgras{Jaume} \textsuperscript{(47)} dormís al luòc d'escotar. Sa
\VUgras{peresa} li mena de \VUgras{reprensions} e quitament de
\VUgras{castigs}. Enfin, \VUgras{Cristian} \textsuperscript{(48)} es tròp
leugièr. Se \VUgras{compòrta} mal e fa pas cap d'esfòrç per s'amendar.

%%K t01-14-1 p
14. --- Sos vesins, al contrari, se compòrtan plan. \VUgras{Ernèst}
\textsuperscript{(51)} aima l'òrdre e la regularitat; se sèrv totjorn de sa
\VUgras{règla} \textsuperscript{(50)} e de son \VUgras{lapis}
\textsuperscript{(49)}. Es \VUgras{extèrne}.

%%K t01-14-2 p
\VUgras{Francés} \textsuperscript{(52)} es \VUgras{intèrne}; pòrta
l'\VUgras{unifòrme} de l'institut; i manja e i dormís. Es un bon
\VUgras{companh}.

%%K t01-14-3 p
Maurici agusa son \VUgras{lapis} amb son \VUgras{cotelet}
\textsuperscript{(56)}; es fòrça curós.

%%K t01-14-4 p
Pòrta totes sos libres, sa \VUgras{gramatica} \textsuperscript{(55)}, son
\VUgras{quasèrn de nòtas} \textsuperscript{(54)} e quitament un
\VUgras{sotamans} \textsuperscript{(53)}. Sa \VUgras{cartabèla d'escòla}
\textsuperscript{(57)} es per tèrra darrièr el.

%%K t01-14-5 p
Las doas taulas del fons de la classa son ocupadas per de \VUgras{bons escolans}
\textsuperscript{(19)}.

%%K t01-tit-6 sub
\VUsaut{4.6mm}
\VUpk{10.2pt}{40}{Dessenh e musica.}
\VUsaut{3.4mm}

%%K t01-15-1 p
15. --- Dins la \VUgras{sala de dessenh} i a de polits \VUgras{modèls} pendolats
a la paret e una \VUgras{lampa} \textsuperscript{(62)} amb \VUgras{abatjorn},
pendolada al plafon. Lo \VUgras{professor} \textsuperscript{(60)} es fòrça
pacient; corregís los \VUgras{dessenhs} \textsuperscript{(61)} dels escolans e lor
ensenha de dessenhar un \VUgras{bust} \textsuperscript{(59)} d'après natura.

%%K t01-16-1 p
16. --- La \VUgras{sala de musica} sembla fòrça interessanta. Aquí s'apren de
legir la \VUgras{musica}, de solfejar las \VUgras{nòtas de la gama}
\textsuperscript{(67)} escrichas sus la \VUgras{portada} (do, re, mi, fa, sol, la,
si\,=\,C. D. E. F. G. A. B.), de reconéisser las \VUgras{claus} e de cantar de
polidas \VUgras{melodias}. Se pausan las particions sul \VUgras{pupitre}
\textsuperscript{(65)}. Un dels escolans \VUgras{jòga del pianò}
\textsuperscript{(64)} e lo \VUgras{professor de musica} \textsuperscript{(66)}
\VUgras{bat la mesura} \textsuperscript{(68)}.

%%K t01-17-1 p
17. --- Començam de veire un canton de l'\VUgras{obrador} de \VUgras{trabalhs
manuals} \textsuperscript{(69)}, ont los dròlles pòdon aprene de rabotar la fusta
e de limar lo fèrre. N'i a pas dins totes los instituts.

%%K t01-tit-7 sub
\VUsaut{5.0mm}
\VUpk{10.2pt}{40}{La sala de sciéncias.}
\VUsaut{3.6mm}

%%K t01-18-1 p
18. --- Lo professor de matematicas ensenha als escolans la \VUgras{geometria,
l'aritmetica, lo calcul} e lo \VUgras{sistèma metric}. Vesèm sus la taula las
principalas figuras de geometria: un \VUgras{cercle} \textsuperscript{(a)}, un
\VUgras{carrat} \textsuperscript{(b)}, un \VUgras{triangle}
\textsuperscript{(c)}, una \VUgras{linha} \textsuperscript{(d)}.

%%K t01-18-2 p
Vesèm tanben una \VUgras{operacion}; es pas una \VUgras{adicion}, ni una
\VUgras{sostraccion}, ni una \VUgras{multiplicacion}: es una \VUgras{division}
\textsuperscript{(e)}.

%%K t01-19-1 p
19. --- Sus la lamina del sistèma metric destriam: un \VUgras{mètre}
\textsuperscript{(a)}, una \VUgras{balança} \textsuperscript{(b)} e sos
\VUgras{pes}, un \VUgras{litre} \textsuperscript{(e)}, un \VUgras{decalitre}
\textsuperscript{(f)}, un \VUgras{decilitre} e un \VUgras{mètre cubic}
\textsuperscript{(d)}.

%%K t01-19-2 p
Los escolans estúdian tanben las \VUgras{sciéncias naturalas}
\textsuperscript{(11)} --- zoologia \textsuperscript{(a)}, botanica
\textsuperscript{(b)}, geologia \textsuperscript{(c)} --- e l'\VUgras{istòria}
\textsuperscript{(12)}.

%%K t01-19-3 p
Dins l'armari i a los instruments de \VUgras{fisica} \textsuperscript{(15)} e de
\VUgras{quimia} \textsuperscript{(16)}.

%%K t01-19-4 p
Dessús l'armari i a: una \VUgras{esfèra} \textsuperscript{(14)}, un \VUgras{còne}
e un \VUgras{globe terrèstre} \textsuperscript{(13)}.

%%K t01-19-5 p
Dessús las laminas pendola una \VUgras{mapa} que representa las cinc parts del
mond: l'\VUgras{Euròpa} \textsuperscript{(g)}, l'\VUgras{Asia}
\textsuperscript{(e)}, l'\VUgras{Africa} \textsuperscript{(f)},
l'\VUgras{America} \textsuperscript{(ab)} e l'\VUgras{Oceania}
\textsuperscript{(h)}, e los dos grands oceans, lo \VUgras{Pacific}
\textsuperscript{(c)} e l'\VUgras{Atlantic} \textsuperscript{(d)}.
\VUsaut{6.0mm}
\VUfilet{22mm}
